problem:  
Let \(M^k \subset S^n\) be a closed, embedded minimal submanifold of the unit sphere. Denote by \(\lambda_1(M)\) the first nonzero eigenvalue of the Laplace–Beltrami operator on \(M\).  

It is known from Takahashi’s theorem that each coordinate function restricted from \(\mathbb{R}^{n+1}\) is an eigenfunction of \(\Delta_M\) with eigenvalue \(k\), hence \(\lambda_1(M)\le k\), with equality precisely for totally geodesic equators \(S^k \subset S^n\).  

**Open Problem (Spectral Pinching Rigidity for Minimal Submanifolds):**  
Does there exist \(\delta = \delta(k,n) > 0\) such that if a closed minimal submanifold \(M^k \subset S^n\) satisfies 
\[
\lambda_1(M) > k - \delta,
\]
then \(M\) must be congruent to a totally geodesic equator \(S^k \subset S^n\)?  

Equivalently, can a quantitative spectral “gap rigidity” principle be established: among minimal submanifolds of the sphere, those whose first eigenvalue comes within \(\delta(k,n)\) of the maximal value \(k\) must exhibit geometric rigidity forcing them to be standard?  

Why is it a "good" problem:  
This formulation probes the **quantitative rigidity** of minimal submanifolds through spectral data. It lies at the crossroads of minimal submanifold theory, spectral geometry, and global analysis. The problem seeks a “spectral pinching implies geometric rigidity” result, akin to classical curvature pinching theorems but in the context of eigenvalues—connecting analytic invariants (eigenvalues) with intrinsic and extrinsic geometry (minimality, curvature, shape operator).  

It is a *good* problem because:  
- It respects known results (Takahashi’s theorem) while extending them into an unknown quantitative direction.  
- It is conceptually simple—asking whether closeness in spectrum forces geometric closeness to the equator—yet potentially deep, invoking techniques from geometric measure theory, stability analysis, and eigenvalue comparison.  
- It bridges analytic and geometric rigidity phenomena, potentially leading to new geometric insights and pinching-type theorems for minimal immersions.  
Even partial results (such as establishing the existence of such a \(\delta(k,n)\) under curvature or stability hypotheses) would represent meaningful advances in the field, illustrating profound interplay between spectral and differential geometry.